% !PRINTED 511
\weberpage{511}

A unit is an integral functional whose reciprocal is likewise integral, and this reciprocal functional is itself a unit. Any integral functional whatever is divisible by a unit. In general the theorem holds:

9. \gesperrt{The product and the quotient of two units are again units.}

For if $\varepsilon_1, \varepsilon_2$ are two units, then $\varepsilon_1\varepsilon_2$ and $1:\varepsilon_1\varepsilon_2$ are integral functionals, hence $\varepsilon_1\varepsilon_2$ is a unit, and likewise $\varepsilon_1:\varepsilon_2$ and $\varepsilon_2:\varepsilon_1$ are integral.

10. \gesperrt{If an integral functional $\mu$ is divisible by another, $\alpha$, then every functional $\mu'$ associated with $\mu$ is also divisible by every functional $\alpha'$ associated with $\alpha$.}

For if $\mu:\alpha$ is integral and $\varepsilon, \varepsilon_1$ are units, then $\mu\varepsilon:\alpha\varepsilon_1$ is also integral.

11. \gesperrt{If $\alpha$ is associated with $\beta$ and with $\gamma$, then $\beta$ and $\gamma$ are also associated with one another.}

For if $\beta=\alpha\varepsilon$, $\gamma=\alpha\varepsilon_1$, then $\beta=\gamma\varepsilon:\varepsilon_1$, and $\varepsilon:\varepsilon_1$ is, by 8., a unit.

If $\alpha$ is divisible by $\beta$, then the absolute norm of $\alpha$ is divisible by the absolute norm of $\beta$; for from $\alpha=\beta\gamma$ it follows, by \S.\,136, 5.:

\weq{1}{N_a(\alpha) = N_a(\beta)\, N_a(\gamma).}

From this it results further that the absolute norm of a unit must be $=1$. But the converse also holds:

12. \gesperrt{An integral functional whose absolute norm is $=1$ is a unit.}

For if $\varepsilon$ is an integral functional with the absolute norm 1, then $N(\varepsilon)$ is an integral rational functional with the absolute value 1, and therefore $1:N(\varepsilon)$ is also an integral functional. If we then set
\[
N(\varepsilon) = \varepsilon\varepsilon',
\]
then $\varepsilon'$, as a product of integral functionals (the conjugates of $\varepsilon$), is itself integral, and consequently
\[
\frac{1}{\varepsilon} = \frac{\varepsilon'}{N(\varepsilon)}
\]
is also an integral functional, hence $\varepsilon$ is a unit.
